\documentclass[10pt]{article}

\usepackage[margin=0.82in]{geometry}
\usepackage{amsmath,amssymb}
\usepackage{booktabs,longtable,tabularx,array}
\usepackage[round]{natbib}
\usepackage{microtype}
\usepackage[hidelinks]{hyperref}
\usepackage{xcolor}

\newcommand{\one}{\mathbf{1}}
\newcommand{\pending}[1]{\textcolor{red!70!black}{[PENDING: #1]}}
\newcolumntype{Y}{>{\raggedright\arraybackslash}X}

\title{A standard bequest block:\\
diagnosis of the estate-target failure and the M4 recalibration}
\author{Fertility--Housing project}
\date{16 July 2026}

\begin{document}
\maketitle

\section*{Executive conclusion}

The July-15 attempt to estimate all three bequest parameters against three
PSID estate moments failed for a reason that is now fully diagnosed: two of
the three moments are not objects this model class can or should match. The
dispersion ratio (p90/p50 of estate to income) requires an upper tail of
financial wealth that no parameterization reaches --- and in the data that
tail is business, IRA, and other nonhousing wealth with no model counterpart.
The family-size gap is statistically zero in the PSID (point estimate
$0.101$, bootstrap standard error $0.563$) and is a knife-edge cancellation
across marital strata that a model without marriage cannot represent. Because
the dispersion moment carried an inverse-bootstrap-variance weight, it alone
contributed 84\% of the SMM loss, and chasing it degraded the eleven
established moments by a factor of four relative to the July-15 mortality-only
(M1) benchmark.

The repair is to return to the standard package used by the structural
savings and housing-retirement literatures: a child-blind luxury warm glow
\citep{denardi2004,lockwood2018,nakajima2017}, with its remaining strength and
luxury-shift parameters estimated internally against \emph{level} targets for
late-life wealth --- the median estate-to-income ratio and the median
nonhousing wealth-to-income ratio. The M4 recalibration implements exactly
this. Headline result: \pending{M4 winner loss, estate median fit, nonhousing
median fit, theta0/theta1 estimates, zero-seed comparison, and identification
verdict}.

\section{Diagnosis: what the estate targets did}

A matched PSID/model decomposition of the age 76--84 estate distribution
(both normalized by each dataset's median annual income at those ages) shows
the failure is composition, not curvature:

\begin{center}
\begin{tabular}{lrrrrrr}
\toprule
& \multicolumn{3}{c}{PSID} & \multicolumn{3}{c}{Model (M3 winner)} \\
\cmidrule(lr){2-4}\cmidrule(lr){5-7}
Ages 76--84 & p50 & p75 & p90 & p50 & p75 & p90 \\
\midrule
Total estate      & 6.87 & 17.10 & 36.33 & 6.27 & 8.87 & 11.79 \\
Nonhousing wealth & 2.37 & 10.34 & 26.27 & 0.00 & 0.00 & 0.00 \\
Home equity       & 3.19 &  6.40 & 11.61 & 6.27 & 8.87 & 11.79 \\
\bottomrule
\end{tabular}
\end{center}

The model matched the estate median almost exactly --- with the wrong balance
sheet. Housing supplies 99.9\% of top-decile estate wealth in the model
versus 18.6\% in the PSID, and every old model household carries collateralized
debt (liquid position $b<0$ through the 90th percentile). The mechanism is
transparent: after age 62 the model's only borrowing is mortgage debt with a
constant floor $b \geq -\phi p H$ and no amortization; the pension is a
common scalar, so no late-life risk motivates liquid saving; and the warm
glow values housing at gross price while a living sale nets $(1-\psi)pH$, so
dying in place dominates. Home equity is therefore the cheapest estate
vehicle, and the model uses only it.

Critically, this is a property of the July-15 \emph{calibration point}, not
of the model. The M1 mortality-only winner --- same constraints, no estimated
bequest parameters --- delivered a positive old nonhousing median of $2.00$
against the corrected reference-person PSID target of $1.91$, while fitting the twelve shared non-estate moments
at a summed loss of $3.84$; the M3 winner fits the same twelve at $16.92$.
The unreachable dispersion row purchased nothing and paid with everything
else, including flipping the young renter liquid-wealth moment to the wrong
sign ($-0.18$ versus a target of $+0.18$, which M1 matched almost exactly).

\section{The standard specification}

M4 uses the bequest utility
\begin{equation}
B(W) \;=\; \theta_0\,
\frac{(\theta_1+W)^{1-\sigma}-\theta_1^{1-\sigma}}{1-\sigma},
\qquad W=\max\{b+pH,\,0\},\quad \sigma=2,
\label{eq:m4bequest}
\end{equation}
where $W$ is the household's net worth at death: the liquid position (with
mortgage debt as negative $b$) plus gross housing value, floored at zero.
This is the luxury warm glow of \citet{denardi2004}; \citet{lockwood2018}
estimates the same form, and \citet{nakajima2017} estimate it inside a
housing-retirement model, which makes their calibration the closest anchor
for this project.

\paragraph{Internally estimated shifter and external child restriction.}
Both $\theta_0$ and $\theta_1$ are estimated internally. The value $0.25$ is
only one optimizer start, motivated by the published \citet{nakajima2017}
estimates (Table~I, Panel~B: $\gamma=20.534$ and $\zeta=\$7{,}619$ per year
in 2000 dollars) and their retired-homeowner income scale; it is not imposed
as a parameter in our normalization. Dispersed starts span the predeclared
log domain $\theta_1\in[0.02,16]$, which also covers the much larger shift in
\citet[eq.~(3), Table~3]{denardifrenchjones2010} ($k=\$273{,}000$ in 1998
dollars). \citet[eq.~(8), Table~4]{denardi2004} report $\phi_2=11.6$ in model
units without a dollar mapping, so that value does not supply a conversion.
The child multiplier alone is externally restricted, $\theta_n=0$. The
post-calibration Jacobian and profile must show whether the two wealth levels
actually separate $\theta_1$ from $\theta_0$, discounting, and the housing
block; otherwise M4 is reported as weakly identified rather than repaired by
fixing $\theta_1$ after the fact.

\paragraph{Why a child-blind motive in a paper about children.} Children in
this model direct resources through housing while parents are alive: they
raise required space, parents buy larger owner-occupied houses, transaction
costs make those houses sticky into old age, and housing is the dominant
component of estates. Parents therefore leave larger estates without any
child-directed taste for posthumous giving. That is where the evidence points.
In \citet[Tables~4--5]{hurd1989}, households with children do not retain more
wealth than observationally comparable households without children.
\citet[pp.~208--209 and Tables~3--4]{kopczuk2007} reject children as a
deterministic indicator of the motive and estimate an egoistic strength common
to households for whom the motive is operative. \citet[pp.~750, 789]{denardi2025}
instead organize luxury bequest saving by marital status, permanent income,
and age. This does not prove that children never matter:
\citet{kvaerner2023} supplies important health-shock evidence for intentional
bequests, and inter vivos transfers can be child-directed
\citep{mcgarry1999}. It does support the parsimonious restriction that the
terminal warm glow itself does not scale mechanically with child count. Our
own sample reinforces that restriction: the family-size gap in median estates
at ages 65--75 is $0.101$ with a bootstrap standard error of $0.563$.

\section{The M4 calibration}

Fourteen moments, thirteen estimated parameters (the eleven clean-frontier
coordinates plus $\theta_0$ and $\theta_1$). The twelve established fertility, ownership,
rooms, and young-liquid-wealth moments are unchanged. The two late-life
wealth targets are levels, in the style of the median asset profiles targeted
by \citet{denardifrenchjones2010} and \citet{nakajima2017}:

\begin{center}
\begin{tabular}{lrl}
\toprule
Moment & Target & Disciplines \\
\midrule
Median total estate wealth / income, ages 76--84 & 6.501 & $(\theta_0,\theta_1)$ jointly \\
Median nonhousing wealth / income, ages 65--75 & 1.908 & $(\theta_0,\theta_1)$ and composition \\
\bottomrule
\end{tabular}
\end{center}

Weights are inverse person-bootstrap variances (499 replications, person-level
clustering). The estate dispersion ratio, the family-size estate gap, the
housing/nonhousing decomposition, and the ownership path by age are reported
as untargeted diagnostics. The search enforces nested-model dominance: the
exact $\theta_0=0$ submodel (the M1 winner under the M4 objective) is injected
as a seed, and the run fails if the free search ends above it.

\section{Results}

The initial six-chain winner had loss $12.654$, but the predeclared conditional
$\theta_1$ profile found a lower basin. A full 13-parameter polish from the
profile minimum produced the final point below: two bit-identical tight
re-solves give a 14-moment loss of $4.402$. Table~\ref{tab:m4fit} shows the
complete fit. The estate median misses by only $0.003$, the nonhousing median
misses by $0.095$, and the established twelve-moment loss is $3.650$, below
the predeclared $4.42$ ceiling. M4 therefore passes every numerical fit gate.

\begin{longtable}{p{0.38\textwidth}rrrrr}
\caption{M4 target fit}\label{tab:m4fit}\\
\toprule
Moment & Target & Model & Gap & Weight & Loss \\
\midrule
\endfirsthead
\toprule
Moment & Target & Model & Gap & Weight & Loss \\
\midrule
\endhead
TFR & 1.918 & 2.000 & 0.082 & 20.000 & 0.134 \\
Childless rate & 0.188 & 0.188 & $-0.000$ & 20.000 & 0.000 \\
Ownership rate & 0.575 & 0.715 & 0.139 & 100.000 & 1.944 \\
Parent--nonparent ownership gap & 0.168 & 0.207 & 0.040 & 45.000 & 0.070 \\
Housing increment, zero to one child & 0.664 & 0.759 & 0.095 & 14.000 & 0.126 \\
Childless renter mean rooms & 3.805 & 3.710 & $-0.095$ & 6.000 & 0.054 \\
Childless owner share in 6+ rooms & 0.596 & 0.724 & 0.128 & 25.000 & 0.409 \\
Young childless renter liquid wealth/income & 0.179 & 0.259 & 0.079 & 12.000 & 0.075 \\
Childless owner--renter room gap & 2.419 & 2.279 & $-0.139$ & 12.000 & 0.233 \\
Ownership rate, ages 25--34 & 0.341 & 0.344 & 0.003 & 80.000 & 0.001 \\
3+ versus 1--2 child room gap & 0.368 & 0.493 & 0.125 & 8.000 & 0.126 \\
Estate wealth/income median, ages 76--84 & 6.501 & 6.505 & 0.003 & 18.586 & 0.000 \\
Nonhousing wealth/income median, ages 65--75 & 1.908 & 2.003 & 0.095 & 83.749 & 0.752 \\
Mean occupied rooms, ages 18--85 & 5.780 & 5.498 & $-0.282$ & 6.000 & 0.477 \\
\midrule
Total & & & & & 4.402 \\
\bottomrule
\end{longtable}

The complete parameter vector is reported in Table~\ref{tab:m4parameters}.
Neither bequest parameter is near a search bound, so the failure is not a
mechanical consequence of the predeclared domains.

\begin{table}[!htb]
\centering
\caption{Estimated M4 parameters and search bounds.}\label{tab:m4parameters}
\begin{tabular}{lrrrl}
\toprule
Parameter & Estimate & Lower & Upper & Near bound \\
\midrule
$\beta$ (annual) & 0.99252 & 0.800 & 0.9995 & No \\
$\alpha_{cons}$ & 0.62122 & 0.020 & 0.980 & No \\
$\bar c_0$ & 1.13244 & 0.000 & 2.000 & No \\
$\bar c_n$ & 0.47052 & 0.000 & 3.000 & No \\
$\bar h_0$ & 0.44449 & 0.050 & 5.800 & No \\
$\bar h_{jump}$ & 1.69570 & 0.000 & 8.000 & No \\
$\bar h_n$ & 0.17997 & 0.000 & 5.000 & No \\
$\psi_{child}$ & 0.22392 & $-3.000$ & 3.000 & No \\
$\kappa_{fert}$ & 2.14745 & 0.020 & 50.000 & No \\
$\chi$ & 1.12893 & 0.100 & 5.000 & No \\
$H_0$ & 7.85561 & 0.200 & 80.000 & No \\
$\theta_0$ & 0.24750 & 0.000 & 8.000 & No \\
$\theta_1$ & 0.42171 & 0.020 & 16.000 & No \\
\bottomrule
\end{tabular}
\end{table}

The external restrictions are $\theta_n=0$ and
$\texttt{tenure\_choice\_kappa}=0$. SSA post-retirement survival is active,
the bequest utility is normalized to zero at a zero estate, and the owner-LTV
taper is inactive.

The untargeted diagnostics show that M4 raises the estate level through the
same housing-heavy mechanism as M3. At ages 76--84, model total estate divided
by the common old-age income level is $6.057$, $8.656$, and $11.415$ at the
50th, 75th, and 90th percentiles; the matched PSID levels are $6.867$,
$17.101$, and $36.335$. Model nonhousing wealth is zero through the 90th
percentile, so model housing equity equals total estate at all three
quantiles. The top-estate-decile housing share is $99.9\%$ in M4 versus
$18.6\%$ in the PSID, and the estate--income correlation is zero rather than
$0.393$. Thus the model can move the estate median but still cannot generate
the financial-wealth tail.

The model estate p90/p50 ratio is $1.884$ against $3.448$ in the PSID. Its
untargeted 2+-minus-1-child median estate gap is $4.383$, versus a data estimate
of $0.101$ with standard error $0.563$. The same mechanism pushes the ownership
rate to $0.983$ at age 62 and keeps it at $0.970$ at age 78, far above the ACS
path displayed in the standard lifecycle figure. Table~\ref{tab:arms} locates
M4 between the clean M1 benchmark and the rejected M3 estate-tail exercise.

\begin{table}[!htb]
\centering
\caption{Comparison with the preceding bequest specifications}\label{tab:arms}
\begin{tabular}{lrrr}
\toprule
 & M1 & M3 & M4 \\
\midrule
Comparable 12-moment loss & 3.841 & 16.921 & 8.242 \\
Aggregate ownership rate & 0.607 & 0.713 & 0.782 \\
Old-age ownership rate & 0.901 & 0.977 & 0.983 \\
Estate median, ages 76--84 & --- & 6.267 & 6.057 \\
Nonhousing median, ages 65--75 & 2.003 & $-1.245$ & 2.003 \\
\bottomrule
\end{tabular}
\end{table}

\pending{Nested-zero and identification verdicts: tight loss at the estimated
$\theta_0$ versus the $\theta_0=0$ seed; Jacobian/profile evidence on whether
$\theta_1$ is separately identified; interpretation of whether the wealth
levels support a positive bequest strength and determinate luxury shift.}

\section{What remains open}

The model still cannot produce the PSID estate tail, by design: that tail is
financial and business wealth held by households with resources far above the
pension-financed median, and matching it would require an earnings or
entrepreneurial tail in the style of \citet{denardi2004} or medical-expense
risk in the style of \citet{denardifrenchjones2010}. Neither belongs in this
paper unless upper-tail wealth becomes substantive for the housing--fertility
question. The old-age leverage margin (no amortization of owner debt) also
remains; policy counterfactuals that operate through late-life home equity
should not be run off this block without further work.

\small
\begin{thebibliography}{99}

\bibitem[De Nardi(2004)]{denardi2004}
De Nardi, M. (2004).
Wealth inequality and intergenerational links.
\emph{Review of Economic Studies} 71(3), 743--768.

\bibitem[De Nardi et al.(2010)]{denardifrenchjones2010}
De Nardi, M., E.~French, and J.~B. Jones (2010).
Why do the elderly save? The role of medical expenses.
\emph{Journal of Political Economy} 118(1), 39--75.

\bibitem[De Nardi et al.(2025)]{denardi2025}
De Nardi, M., E.~French, J.~B. Jones, and R.~McGee (2025).
Why do couples and singles save during retirement? Household heterogeneity
and its aggregate implications.
\emph{Journal of Political Economy} 133(3), 750--792.

\bibitem[Hurd(1989)]{hurd1989}
Hurd, M.~D. (1989).
Mortality risk and bequests.
\emph{Econometrica} 57(4), 779--813.

\bibitem[Kopczuk and Lupton(2007)]{kopczuk2007}
Kopczuk, W. and J.~P. Lupton (2007).
To leave or not to leave: The distribution of bequest motives.
\emph{Review of Economic Studies} 74(1), 207--235.

\bibitem[Kv\ae rner(2023)]{kvaerner2023}
Kv\ae rner, J. (2023).
How large are bequest motives? Estimates based on health shocks.
\emph{Review of Financial Studies} 36(8), 3382--3422.

\bibitem[Lockwood(2018)]{lockwood2018}
Lockwood, L.~M. (2018).
Incidental bequests and the choice to self-insure late-life risks.
\emph{American Economic Review} 108(9), 2513--2550.

\bibitem[McGarry(1999)]{mcgarry1999}
McGarry, K. (1999).
Inter vivos transfers and intended bequests.
\emph{Journal of Public Economics} 73(3), 321--351.

\bibitem[Nakajima and Telyukova(2017)]{nakajima2017}
Nakajima, M. and I.~A. Telyukova (2017).
Reverse mortgage loans: A quantitative analysis.
\emph{Journal of Finance} 72(2), 911--950.

\end{thebibliography}

\end{document}
